% 02_related.tex — owner: A2 (Round 4). All keys verified present in docs/refs.bib.
\section{Related Work}
\label{sec:related}

\subsection{LLM-based vulnerability detection}
PrimeVul~\cite{ding2025primevul} re-anchored function-level vulnerability
detection on a deduplicated, chronologically split C/C++ corpus with paired
vulnerable/patched functions, showing large optimism gaps in legacy benchmarks
and proposing strict paired metrics (including detection at constrained false
positive rate). Gao et al.~\cite{han2023howfar} reached similar conclusions
for LLMs specifically, and LLM4Vuln~\cite{sun2024llm4vuln} decouples prompt,
knowledge, and reasoning factors --- our E0 arms follow its controlled-prompt
logic. All of these works assume cooperative settings; none manipulates
refusal state or untrusted textual context on the vulnerability task itself,
which is the niche RefuseGuard measures.

\subsection{Over-refusal and defensive refusal}
XSTest~\cite{rottger2024xstest} and OR-Bench~\cite{cui2025orbench} established
over-refusal as a measurable failure mode and supply our calibration and
contrast corpora. Campbell et al.~\cite{campbell2026defensive} report that
cyber-defence requests with security-sensitive wording are refused about
$2.7\times$ more than neutral equivalents; TabooRAG~\cite{li2026taboorag}
shows alignment can be exploited to block RAG answers; and
\cite{li2026beyondrefusal} studies vulnerability-analysis utility of aligned
versus safety-ablated models. Mitigation work spans query-side gating
(SCANS~\cite{cao2025scans}), activation steering
(ACTOR~\cite{dabas2025actor}), and dual-objective alignment
(DOOR~\cite{zhao2025door}). These studies target generic chat safety; none
evaluates mitigation \emph{on a security-analysis task with verifiable
correctness}, and none keeps the unsafe-compliance side of the ledger ---
exactly the joint evaluation our E8 performs. Our reproduction gate exists
precisely because \cite{campbell2026defensive}'s effect is measured on
different tasks (CTF-style defence prompts), and we show it does not transfer
to vulnerability analysis for the models we test.

\subsection{Indirect prompt injection in code contexts}
BIPIA~\cite{yi2023bipia} benchmarked IPI through externally sourced content
and separated informational from instructional text; Liu et
al.~\cite{liu2024autopi} demonstrated optimized universal injection suffixes;
CodeSentinel~\cite{cheng2026codesentinel} defends code contexts with a
three-layer sanitizer targeting comments, strings, and identifiers. Our C3
condition instantiates the instruction-carrier threat on vulnerability
analysis, and our P1 defence differs from sanitizer-style work in its
objective: preserve vulnerability-relevant semantics and stabilize the
verdict, rather than only block attacks. Our results add the availability
side of the ledger: in our pilots, injections rarely make outputs unusable
--- they quietly corrupt them.

\subsection{Safety--utility trade-offs}
Fine-tuning can silently erode safety alignment~\cite{qi2024finetuning}, and
cyber-domain fine-tuning trades safety against task
performance~\cite{cyberllminstruct2025}. Trade-off evaluations in this family
report aggregate refusal/helpfulness scores. RefuseGuard instead couples both
sides to \emph{verifiable task correctness} (vulnerable/benign verdicts, CWE,
localization) and to an inference-time defence with fallback, and treats
``fixing'' over-refusal by making the model indiscriminately compliant as a
first-class failure mode --- the refusal-suppression phenomenon we document in
Section~\ref{sec:results} (RQ6).

\subsection{Positioning}
We claim neither the observation of defensive refusal
(\cite{campbell2026defensive} owns it) nor the first IPI-in-code defence
(\cite{cheng2026codesentinel,yi2023bipia} own that). Our contribution is the
first (to our verified knowledge) \emph{joint} refusal--utility--safety
evaluation of LLM-based vulnerability analysis under paired clean vs.\
security-charged code contexts, together with a reproduction-gate methodology
that converted an expected phenomenon into an honest negative result.
